% ═══════════════════════════════════════════════════════════════
% AB-01b: Sich vorstellen + Verb ser
% Abgeleitet aus LE-01b (Score 9.2/10)
% ═══════════════════════════════════════════════════════════════

\documentclass[11pt, a4paper]{article}

% ═══════════════════════════════════════════════════════════════
% SW-MODUS TOGGLE
% ═══════════════════════════════════════════════════════════════
% Setze \swmodustrue für Schwarz-Weiß-Druck
\newif\ifswmodus
\swmodusfalse    % Standard: Farbe

\usepackage[utf8]{inputenc}
\usepackage[T1]{fontenc}
\usepackage[ngerman]{babel}
\usepackage{helvet}
\renewcommand{\familydefault}{\sfdefault}

\usepackage[margin=2cm]{geometry}
\usepackage{enumitem}
\usepackage{tabularx}
\usepackage{booktabs}

\usepackage{xcolor}
\usepackage{tcolorbox}
\tcbuselibrary{skins}

\usepackage{fancyhdr}
\usepackage{lastpage}
\usepackage{needspace}          % Für \needspace (Seitenumbruch-Kontrolle)

% Farben - Basis
\definecolor{spanisch-color}{HTML}{c41e3a}
\definecolor{grau-color}{HTML}{888888}
\definecolor{hellgrau-color}{HTML}{f5f5f5}
\definecolor{orange-color}{HTML}{b35c00}

% SW-Farben (druckerfreundlich: keine Füllflächen, nur Linien)
\definecolor{sw-dark}{HTML}{000000}
\definecolor{sw-gray}{HTML}{666666}
\definecolor{sw-white}{HTML}{ffffff}

% Bedingte Farbzuweisung
\ifswmodus
    \colorlet{spanisch}{sw-dark}       % Rahmen/Text: schwarz
    \colorlet{grau}{sw-gray}           % Sekundärtext: dunkelgrau
    \colorlet{hellgrau}{sw-white}      % Hintergrund: weiß (keine Füllung)
    \colorlet{orange}{sw-dark}         % Hinweis-Rahmen: schwarz
\else
    \colorlet{spanisch}{spanisch-color}
    \colorlet{grau}{grau-color}
    \colorlet{hellgrau}{hellgrau-color}
    \colorlet{orange}{orange-color}
\fi

\newcommand{\fachfarbe}{spanisch}

% Variablen
\newcommand{\thema}{Sich vorstellen + Verb ser}
\newcommand{\sequenz}{01}
\newcommand{\block}{b}
\newcommand{\fach}{Spanisch}
\newcommand{\klasse}{7}
\newcommand{\maxpunkte}{20}

% Kopf-/Fußzeile
\pagestyle{fancy}
\fancyhf{}
\renewcommand{\headrulewidth}{0.4pt}
\renewcommand{\footrulewidth}{0.4pt}
\lhead{\fach{} -- Klasse \klasse}
\chead{AB-\sequenz\block}
\rhead{}
\lfoot{}
\rfoot{Seite \thepage\ von \pageref{LastPage}}

% Befehle
\newcommand{\punktebox}[1]{\hfill\fbox{\small \_/#1}}
\newcommand{\luecke}[1]{\rule{#1}{0.4pt}}
\newcommand{\es}[1]{\textcolor{\fachfarbe}{\textbf{#1}}}

% Merkkasten (SW-kompatibel: kein gefüllter Titelbalken)
\newtcolorbox{merkkasten}[1][]{
    colback=hellgrau,
    colframe=\fachfarbe,
    colbacktitle=white,
    coltitle=\fachfarbe,
    boxrule=1pt,
    arc=0pt,
    left=8pt, right=8pt, top=8pt, bottom=8pt,
    fonttitle=\bfseries,
    attach boxed title to top left={yshift=-2mm, xshift=5mm},
    boxed title style={boxrule=0pt, colback=white},
    title=#1
}

% Hinweis (SW-kompatibel: kein gefüllter Titelbalken)
\newtcolorbox{hinweis}{
    colback=white,
    colframe=orange,
    colbacktitle=white,
    coltitle=orange,
    boxrule=1pt,
    arc=0pt,
    left=5pt, right=5pt, top=5pt, bottom=5pt,
    fonttitle=\bfseries,
    attach boxed title to top left={yshift=-2mm, xshift=5mm},
    boxed title style={boxrule=0pt, colback=white},
    title=Hinweis
}

% Aufgabe (kein Seitenumbruch innerhalb)
\newenvironment{aufgabe}[1]{%
    \needspace{5\baselineskip}%     % Mindestens 5 Zeilen Platz, sonst neue Seite
    \nopagebreak[4]%
    \vspace{10pt}
    \noindent\textbf{\textcolor{\fachfarbe}{Aufgabe #1}}
    \nopagebreak[4]%
    \vspace{5pt}
    \nopagebreak[4]%
    \begin{list}{}{\leftmargin=0pt}
    \item[]
}{%
    \end{list}
}

% ═══════════════════════════════════════════════════════════════
\begin{document}

% Kopfbereich
\begin{center}
    {\Large\bfseries\textcolor{\fachfarbe}{Arbeitsblatt: \thema}}
\end{center}

\vspace{5pt}

\noindent
\begin{tabularx}{\textwidth}{|X|X|X|}
\hline
\textbf{Name:} \luecke{4cm} &
\textbf{Klasse:} \klasse &
\textbf{Datum:} \luecke{2cm} \\
\hline
\end{tabularx}

\vspace{15pt}

% ═══════════════════════════════════════════════════════════════
% MERKKASTEN: Verb ser ($\rightarrow$ FZ1)
% ═══════════════════════════════════════════════════════════════

\begin{merkkasten}[Merkkasten: Das Verb \textit{ser} (sein)]

\begin{center}
\begin{tabular}{|l|l|l|}
\hline
\textbf{Person} & \textbf{Pronomen} & \textbf{ser} \\
\hline
1. Sg. & yo (ich) & \luecke{2cm} \\
\hline
2. Sg. & tú (du) & \luecke{2cm} \\
\hline
3. Sg. & él/ella (er/sie) & \luecke{2cm} \\
\hline
1. Pl. & nosotros/-as (wir) & \luecke{2cm} \\
\hline
\end{tabular}
\end{center}

\vspace{0.5em}
\textit{Übertrage die Formen aus dem Lernpfad oder von der Tafel!}
\end{merkkasten}

\vspace{10pt}

% ═══════════════════════════════════════════════════════════════
% AUFGABE 1: Konjugation zuordnen ($\rightarrow$ FZ1, AFB I)
% ═══════════════════════════════════════════════════════════════

\begin{aufgabe}{1}
Verbinde die Pronomen mit der richtigen Form von \textit{ser}. \punktebox{4}
\end{aufgabe}

\begin{center}
\begin{tabular}{rcl}
yo & \luecke{3cm} & es \\[0.8em]
tú & \luecke{3cm} & soy \\[0.8em]
él/ella & \luecke{3cm} & somos \\[0.8em]
nosotros & \luecke{3cm} & eres \\
\end{tabular}
\end{center}

\vspace{15pt}

% ═══════════════════════════════════════════════════════════════
% AUFGABE 2: Lückentext ($\rightarrow$ FZ3, AFB II)
% ═══════════════════════════════════════════════════════════════

\begin{aufgabe}{2}
Setze die richtige Form von \textit{ser} ein. \punktebox{4}
\end{aufgabe}

\begin{enumerate}[label=\alph*)]
    \item Me llamo María. \luecke{2cm} de Alemania.
    \item ¿Cómo te llamas? -- \luecke{2cm} Pedro.
    \item Él \luecke{2cm} de España.
    \item Nosotros \luecke{2cm} de Rostock.
\end{enumerate}

\vspace{15pt}

% ═══════════════════════════════════════════════════════════════
% AUFGABE 3: Dialog vervollständigen ($\rightarrow$ FZ2, AFB II)
% ═══════════════════════════════════════════════════════════════

\begin{aufgabe}{3}
Vervollständige den Dialog. Nutze die Redemittel aus dem Kasten. \punktebox{6}
\end{aufgabe}

\begin{hinweis}
\textbf{Redemittel:} ¿Cómo te llamas? | Me llamo... | ¿De dónde eres? | Soy de... | ¿Y tú?
\end{hinweis}

\vspace{0.5em}

\noindent
\textbf{A:} \es{¡Hola!} \luecke{5cm}

\textbf{B:} \luecke{6cm} \es{¿Y tú?}

\textbf{A:} \es{Soy Carlos.} \luecke{4cm}

\textbf{B:} \luecke{6cm}

\textbf{A:} \es{¡Mucho gusto!}

\textbf{B:} \luecke{4cm}

\vspace{15pt}

% ═══════════════════════════════════════════════════════════════
% AUFGABE 4: Freier Dialog ($\rightarrow$ FZ4, AFB III)
% ═══════════════════════════════════════════════════════════════

\begin{aufgabe}{4}
Schreibe einen eigenen Vorstellungsdialog (mind. 6 Zeilen). \punktebox{6}
\end{aufgabe}

\begin{hinweis}
Verwende: Begrüßung, Name, Herkunft, Verabschiedung. Bonus: Erkläre, wann man \textit{tú} und wann \textit{usted} verwendet!
\end{hinweis}

\noindent\rule{\textwidth}{0.4pt}\vspace{8pt}
\noindent\rule{\textwidth}{0.4pt}\vspace{8pt}
\noindent\rule{\textwidth}{0.4pt}\vspace{8pt}
\noindent\rule{\textwidth}{0.4pt}\vspace{8pt}
\noindent\rule{\textwidth}{0.4pt}\vspace{8pt}
\noindent\rule{\textwidth}{0.4pt}\vspace{8pt}

\vspace{20pt}
\hrule
\vspace{10pt}

\begin{flushright}
\textbf{Punkte:} \luecke{1cm} / \maxpunkte
\end{flushright}

\end{document}
